% Every LaTeXML box's width reaches the HTML as the CSS custom property `--ltx-box-width` (DG 2026-11-14, "Move them
% all now"): a minipage, a \parbox, a \makebox, a transformed box (outer and inner), a wrapfigure, a section lifted out
% of a minipage (QA24); the XML keeps its `width` attribute. A rule's width is its ink and MathML's `mspace` width is
% MathML's own: both stay `width:`; varwidth's own `width:max-content` follows the property and outranks it in LaTeXML.css.
% witness: synthetic (Bruce Q&A probe Q2.4); 2512.24470 (a \section inside a minipage).
% status:  GREEN (QAB57b)
% oracle:  pdflatex 0 errors: each box set at its width.
% Perl:    `style="width:<w>"` from LaTeXML-common.xsl's base-styling (a lifted section: no width).
% ours:    before QAB57b `style="width:<w>"`.
% expect:  0 errors, 0 warnings; each box `style="--ltx-box-width:<w>;…"`, the rule and the mspace `width:`.
\documentclass{article}
\usepackage{graphicx,varwidth,wrapfig}
\begin{document}
\begin{minipage}[t]{3cm}A minipage.\end{minipage} \parbox{2cm}{A parbox.} \makebox[4cm][r]{boxed}

\rotatebox{90}{\parbox{2cm}{Turned}} \rule{3cm}{2pt} $x\hspace{1cm}y$

\begin{varwidth}{8cm}short\end{varwidth}

\begin{wrapfigure}{r}{0.3\textwidth}
\centering W
\end{wrapfigure}
Text beside the wrapfigure.

\begin{minipage}{4cm}
\subsection*{Lifted}
Text in a lifted section.
\end{minipage}
\end{document}
